% =============================
% Knowledge Compilation Map - Conceptual Definitions
% Auto-generated from database.json
% Generated: 2026-04-08T00:35:09.462Z
%
% EDITING INSTRUCTIONS:
% - Definition IDs are preserved in \label{...}. Do NOT edit.
% - Titles (\textbf{...}) are EDITABLE.
% - Statement content (after the title line) is EDITABLE.
% - To sync back to JSON, run: npx tsx scripts/latex-bijection.ts --to-json
% =============================
\documentclass[11pt]{article}

% -------- Packages --------
\usepackage[margin=1in]{geometry}
\usepackage{amsmath, amssymb, amsthm}
\usepackage{mathtools}
\usepackage{enumitem}
\usepackage{hyperref}
\usepackage{cleveref}
\usepackage{xcolor}
\usepackage{natbib}

% -------- Hyperref setup --------
\hypersetup{
  colorlinks=true,
  linkcolor=blue,
  citecolor=blue,
  urlcolor=blue
}

% -------- Theorem styles --------
\theoremstyle{definition}
\newtheorem{definition}{Definition}

% -------- Handy macros --------
\newcommand{\R}{\mathbb{R}}
\newcommand{\N}{\mathbb{N}}
\newcommand{\eps}{\varepsilon}
% Cross-reference commands (rendered as links in the web UI)
\newcommand{\langref}[1]{\textbf{#1}}
\newcommand{\edgeref}[2]{#1 compiles to #2}
\newcommand{\nedgeref}[2]{#1 cannot compile to #2}
\newcommand{\opref}[2]{#1 supports #2}
\newcommand{\nopref}[2]{#2 is unsupported by #1}

% -------- Title info --------
\title{Knowledge Compilation Map: Conceptual Definitions}
\date{\today}

\begin{document}
\maketitle

\begin{definition}\label{kdef:boolean-formula}
\textbf{Boolean Formula} \\
TODO: Define the propositional objects the map ranges over: Boolean formulas as syntactic representations of Boolean functions, together with the usual semantic equivalence relation. \citet{Darwiche_2002}
\end{definition}

\begin{definition}\label{kdef:representation-language}
\textbf{Representation Language} \\
TODO: A representation language is a set of formulas together with an intended semantics and a size measure that makes compilation statements meaningful. \citet{Darwiche_2002}
\end{definition}

\begin{definition}\label{kdef:language-family}
\textbf{Language Family} \\
TODO: A language family is a parameterized collection of representation languages, such as OBDD$_<$ or cSDD$_T$, where the parameter indexes a particular sublanguage. \citet{Darwiche_2002,Amarilli_2018}
\end{definition}

\begin{definition}\label{kdef:succinctness-language}
\textbf{Succinctness Between Languages} \\
TODO: For languages $L_1$ and $L_2$, we write $L_1 \leq_p L_2$ when every formula in $L_1$ can be compiled into $L_2$ in polynomial time. This is the basic notion behind the edge labels in the map. \citet{Darwiche_2002}
\end{definition}

\begin{definition}\label{kdef:succinctness-family}
\textbf{Succinctness Between Families} \\
TODO: For families $F_1$ and $F_2$, one natural reading is $F_1 \leq_p F_2$ when for every language $L_2 \in F_2$ there exists a language $L_1 \in F_1$ such that $L_1 \leq_p L_2$. We will tune the quantifier order carefully before formalizing it in LaTeX. \citet{Darwiche_2002}
\end{definition}

\begin{definition}\label{kdef:tractability}
\textbf{Tractability} \\
TODO: An operation is tractable on a language when it can be performed in polynomial time on every instance represented in that language. \citet{Darwiche_2002}
\end{definition}

\begin{definition}\label{kdef:query-operation}
\textbf{Query Operation} \\
TODO: Query operations are semantic decision tasks such as consistency, validity, entailment, satisfiability, and model counting. \citet{Darwiche_2002}
\end{definition}

\begin{definition}\label{kdef:transformation-operation}
\textbf{Transformation Operation} \\
TODO: Transformation operations are compilation-preserving map-like tasks such as conditioning, forgetting, conjunction, and disjunction. \citet{Darwiche_2002}
\end{definition}

% =============================
% Bibliography
% =============================
\bibliographystyle{plainnat}
\bibliography{refs}

\end{document}
